% Regression: p3_multibond2 — imported current-behavior fixture from the handoff corpus.
% Formula: A_1 (X on bond 1; bonds 2..k) A_2 = B_1 (same) B_2   for all X
% p3_multibond2 — arXiv:1804.04964, eq:lem_inj_eq_ten_2 (lines 1069-1090):
% two injective tensors joined by SEVERAL bonds (a multi-edge), a matrix
% X on the first bond, and the paper's \vdots eliding the remaining bonds.
%
%   A_1 (X on bond 1; bonds 2..k) A_2 = B_1 (same) B_2   for all X
%     ==> A_1 = lambda B_1, A_2 = lambda^{-1} B_2.
%
% Paper ink: two beads with physical legs up, three bent parallel bonds,
% X on the top bond, \vdots between the lower bonds.  This attempt PROBES
% the elision spelling: \tndots placed in the LAST wire row's middle cell,
% hoping it reads as "more parallel bonds".  (The prior p3_multibond fixed
% k = 3 with \tnghost and recorded the missing vertical-ellipsis marker;
% this file tests whether \tndots in a wire row is an acceptable stand-in.)
\documentclass[varwidth=19cm, border=6pt]{standalone}
\usepackage{amsmath, amssymb}
\usepackage{tenkz}
\begin{document}
% A_1 -(X on bond 1, bonds 2..k)- A_2  =  B_1 -(same)- B_2
\[
  \begin{tenkz}[rows={ket,wire,wire}, boundary=none]
    \tn[wires=3, box]{A_1} & \tnX{X}    & \tn[wires=3, box]{A_2} \\
                           & \tnghost{} &                        \\
                           & \tndots    &
  \end{tenkz}
  \;=\;
  \begin{tenkz}[rows={ket,wire,wire}, boundary=none]
    \tn[wires=3, box]{B_1} & \tnX{X}    & \tn[wires=3, box]{B_2} \\
                           & \tnghost{} &                        \\
                           & \tndots    &
  \end{tenkz}
  \ .
\]
\end{document}
